\section{Does the ban bind on trading?}\label{sec:ban}

Before asking what a ban does to prices, which is the subject of the second stage, I first ask whether the rule
changes trading at all. A rule that lets traders only reduce positions should have three visible effects
(\cref{prop:trade}): futures trading should fall, some of the blocked demand should move to the cash market, and every
futures trade should close existing positions. This section tests these effects, looks at what happens at the 95\%
trigger itself, and measures how the October 2025 rules changed the number of bans.

\subsection{Trading on ban days (pre-specified test H8)}
A simple comparison of ban days with other days would mix the effect of the ban with the effect of the heavy trading
that pushed the stock into it. The plan therefore compares each stock's ban days with its own non-ban days at a similar
level of utilisation the day before. Between 80\% and 95\% both kinds of days occur: a stock that crossed 95\% stays in
the ban until it falls to 80\%, while a stock that has not crossed is free. The regression is
\begin{equation}\label{eq:h8}
y_{it}=\beta\,\mathrm{Ban}_{it}+\alpha_{i,b(\omega_{i,t-1})}+\lambda_t+\varepsilon_{it},
\end{equation}
where $\mathrm{Ban}_{it}$ marks ban days, $\alpha_{i,b}$ is a fixed effect for each stock and each 0.05-wide bin of
yesterday's utilisation $\omega_{i,t-1}$, and $\lambda_t$ is a fixed effect for each day, which absorbs market-wide
news. Standard errors are clustered by stock and by day. The sample covers stock-days before October 2025 with
utilisation between 80\% and 95\% the day before, excluding expiry days, when open interest falls mechanically.

In a ban, futures trading falls by 2.0 log points, or about 87\%. Cash-market volume rises by about 7.5\%, and the share of
cash trading that ends in delivery rises by 4.1 percentage points from an average of 26\% (\cref{tab:h8}). All three
effects have the predicted sign. H8 is supported, also after the Holm correction for the four primary hypotheses.

\begin{table}[t]
\caption{In a ban, futures trading collapses and cash trading moves toward delivery}\label{tab:h8}
\exsub{Ban days against the same stock's non-ban days at utilisation of 80\% to 95\% the day before, January 2015 to September 2025}
\small\begin{tabular}{@{}>{\raggedright\arraybackslash\bfseries\color{accent}}p{6.6cm}cc@{}}
\toprule
\thd{Outcome} & \thd{Effect of a ban (s.e.)} & \thd{Mean on non-ban days} \\
\midrule
log(1 + futures contracts traded) & $-2.00$ (0.11) & 8.50 \\
log(1 + cash shares traded) & $+0.073$ (0.027) & 15.80 \\
Delivery share of cash volume & $+0.041$ (0.005) & 0.26 \\
\bottomrule
\end{tabular}

\exsource{Equation \eqref{eq:h8}: stock-by-utilisation-bin and date fixed effects; standard errors clustered by stock and date;
expiry days excluded. 16,270 stock-days, 5,441 of them in a ban, 167 stocks. Source: author's calculations.}
\end{table}

\subsection{Does every trade close a position?}
If the rule binds on every trade, each futures trade during a ban must close a long position on one side and a short
position on the other, so open interest falls by exactly the volume traded. The ratio $R=(\Delta OI+V)/V$ is then zero,
while on a normal day, when most trades open new positions, it is close to one. The plan set the reading in advance: the
identity holds if the mean ratio on ban days is close to zero. Over the whole period it is not: the mean on ban days is
0.42.

The reason became clear once the ratio is plotted over time (\cref{fig:banR}). It changed abruptly in March 2020, at
exactly the point when the enforcement of the rule changed. Until 22 March 2020 violations were judged on end-of-day
positions, so a trader could open and close a position within the day; the median ratio on ban days was then 0.87, close
to a normal day. From 23 March 2020 violations were judged on peak intraday positions, and the median fell to 0.011:
almost every trade closed a position, as the model predicts when the rule binds trade by trade. Under the rule of October
2025, which allows any trade that lowers an entity's end-of-day delta, the median rose again to between 0.19 and
0.36. The fall in futures trading itself was also much larger under intraday checks ($-2.72$ log points) than under
end-of-day checks ($-0.97$; \cref{fig:h8era}). The split by enforcement period was not in the plan; I added it after the
pooled result to find its cause. It shows that the rule changes trading as designed only when positions are checked
during the day.

\begin{figure}[t]
\caption{When intraday checks began in March 2020, every futures trade in a ban started to close positions}\label{fig:banR}
\exsub{Monthly median of $R=(\Delta OI+V)/V$ on ban days, near and next futures, expiry days excluded}
\centering\includegraphics[width=\linewidth]{fig_ban_R.pdf}
\exsource{$R=0$ when every trade closes a long and a short; $R\approx 1$ on normal days. Source: NSE F\&O bhavcopies and ban lists; author's calculations.}
\end{figure}

\begin{figure}[t]
\caption{The ban bites as designed only when positions are checked during the day}\label{fig:h8era}
\exsub{Effect of a ban on log futures contracts (left) and median $R$ on ban days (right), by how violations were judged}
\centering\includegraphics[width=\linewidth]{fig_h8_era.pdf}
\exsource{End-of-day checks to 22 March 2020; peak intraday checks from 23 March 2020 to September 2025; delta rule from October 2025 (from 8 December 2025 shown). Source: author's calculations.}
\end{figure}

\subsection{At the 95\% trigger}
\paragraph{Do traders hold back?} If traders wanted to avoid a ban, utilisation would pile up just below 95\% and thin out
just above it. The model does not predict which way this goes (\cref{sec:model-ban}), so the test is descriptive. Its
population needs care. Among all days on which a stock is not in a ban, utilisation above 95\% is rare by construction,
since a stock that crosses 95\% is in the ban the next day. I therefore use the days on which a stock was below 95\% the
day before and ask how its utilisation is distributed today. The density is smooth through the trigger (\cref{fig:density};
the test of Cattaneo, Jansson and Ma (2020) gives $T=-0.18$, $p=0.86$), and this holds across specifications and
enforcement periods. Utilisation does not stall near the trigger; it speeds up. The share of stocks that cross 95\% the
next day rises steadily from about 7\% at a utilisation of 85\% to about one half at 93\% to 95\%. There is no sign that
traders hold back to avoid a ban.

\begin{figure}[t]
\caption{Utilisation does not bunch below the trigger}\label{fig:density}
\exsub{End-of-day utilisation of stocks that were below 95\% the day before, January 2015 to September 2025}
\centering\includegraphics[width=\linewidth]{fig_density.pdf}
\exsource{Cattaneo, Jansson and Ma test: $T=-0.18$, $p=0.86$. Source: NSE combined open interest files; author's calculations.}
\end{figure}

\paragraph{A sharp break in trading.} The trigger can also be read as a regression discontinuity. For stocks not in a ban,
ending a day above 95\% puts the stock in the ban the next day in 99.7\% of cases, against almost never just below.
Comparing stocks just above and just below the trigger with local linear regressions, futures trading falls by 1.53 log
points, $R$ falls by 0.66 and the delivery share rises by 5.1 percentage points, while cash volume does not jump
(\cref{fig:rdtrading}). These jumps confirm the results of the regression above with a different comparison: stocks that
differ only in whether they ended the previous day just above or just below the trigger.

\begin{figure}[t]
\caption{Futures trading drops sharply at the trigger}\label{fig:rdtrading}
\exsub{Binned means by utilisation at the end of the previous day, stocks not in a ban the day before, expiry days excluded, with linear fits on each side}
\centering\includegraphics[width=\linewidth]{fig_rd_trading.pdf}
\exsource{Estimates in the text are local linear with robust bias-corrected inference (Calonico, Cattaneo and Titiunik, 2014). January 2015 to September 2025. Source: author's calculations.}
\end{figure}

\subsection{How the 2025 rules changed the number of bans}
SEBI's stated aim in redesigning position limits in 2025 was to reduce ``spurious'' bans. Before the reform, bans came
from rising open interest rather than from changes in the limit: on the day utilisation crossed 95\%, the limit had
fallen for only 8 of 1,744 entries, and only 7 of them (0.4\%) would not have crossed at the previous day's limit. The
reform changed two things at once: open interest came to be measured in delta-adjusted terms, and the new formula cut
every stock's limit, by a quarter to a half (\cref{fig:mwplcut}).

\begin{figure}[t]
\caption{The October 2025 formula cut every stock's limit, by a quarter to a half}\label{fig:mwplcut}
\exsub{Ratio of the new to the old MWPL on 1 October 2025, all 208 F\&O stocks}
\centering\includegraphics[width=0.9\linewidth]{fig_mwpl_ratio.pdf}
\exsource{Source: NSE combined open interest files of 30 September and 1 October 2025; author's calculations.}
\end{figure}

To separate the two parts, I simulate the ban rule (enter the day after utilisation exceeds 95\%, leave the day after it is
at or below 80\%) for the stocks in the segment on both 30 September and 1 October 2025. Applied to the actual utilisation
after the reform, the simulation reproduces the official ban lists exactly. I then apply it with open interest measured
gross or in delta terms, and with the new limit or the limit the old formula would have given.\footnote{The old-formula
limit is approximated by scaling each stock's new limit by its ratio of old to new limit on 1 October 2025.} Under the old
rule the same stocks would have spent about 3.7 times as many days in a ban in the first year (\cref{tab:spurious}). The
two parts of the reform pull in opposite directions. Measuring open interest in delta terms removes almost every ban,
because options now count only by their delta and offsetting positions of the same client net out; the lower limit on
its own would have quadrupled them. The combination cut ban days by about three quarters (\cref{fig:spurious}). Whether
the new rules have also changed what a ban does to prices needs more post-reform data, and is part of the second stage.

\begin{table}[t]
\caption{The new rule produced about a quarter of the ban days the old rule would have}\label{tab:spurious}
\exsub{Bans from 1 October 2025 to 27 October 2026 for the stocks in the segment on 30 September and 1 October 2025}
\small\begin{tabular}{@{}>{\bfseries\color{accent}}lccc@{}}
\toprule
\thd{Open interest and limit} & \thd{Ban stock-days} & \thd{Spells} & \thd{Stocks} \\
\midrule
Delta-adjusted, new limit (the actual rule) & 412 & 45 & 11 \\
Gross, old-formula limit (the old rule) & 1,505 & 112 & 36 \\
Gross, new limit & 6,256 & 296 & 85 \\
Delta-adjusted, old-formula limit & 1 & 1 & 1 \\
\bottomrule
\end{tabular}

\exsource{Simulated ban rule, validated against the official lists. The 45 spells include one that began on 30 September 2025. Source: NSE combined open interest files; author's calculations.}
\end{table}

\begin{figure}[t]
\caption{Measuring open interest in delta terms removes almost all bans; the lower limit alone would have multiplied them}\label{fig:spurious}
\exsub{Ban stock-days, 1 October 2025 to 27 October 2026, under the actual and three counterfactual rules}
\centering\includegraphics[width=\linewidth]{fig_spurious.pdf}
\exsource{Source: as for \cref{tab:spurious}.}
\end{figure}
